\documentclass[a4paper]{article}
\usepackage{amsfonts}
\usepackage{amsmath}
\usepackage{amssymb}
\usepackage{graphicx}
\begin{document}
\noindent \large Auteur: Siebe Haché \\
\noindent \large Studierichting: Informatica\\
\noindent \large Oefeningenreeks 1C nr. 31.15\\

\newcommand{\Bgtan}{\operatorname{Bgtan}}
\newcommand{\cis}{\operatorname{cis}}

\medskip
\normalsize
Bereken in $\mathbb{C}$ door gebruik van goniometrische vormen: $\left(\dfrac{-3-5i}{4+i}\right)^5$ \\

$\dfrac{-3-5i}{4+i} = \dfrac{(-3-5i)(4-i)}{(4+i)(4-i)} = \dfrac{-12+3i-20i-5}{16+1} = \dfrac{-17-17i}{17} = -1-i$ \\

$z = -1-i$ \\

$r = \sqrt{(-1)^2+(-1)^2} = \sqrt{2}$ \\

$\theta = \pi + \Bgtan\left(\dfrac{-1}{-1}\right) = \pi + \dfrac{\pi}{4} = \dfrac{5\pi}{4}$ \\

$z = \sqrt{2} \cis \left(\dfrac{5\pi}{4}\right)$ \\

$z^5 = \left(\sqrt{2}\right)^5 \cis \left(5 \cdot \dfrac{5\pi}{4}\right) = 4\sqrt{2} \cis \left(\dfrac{25\pi}{4}\right)$ \\

$z^5 = 4\sqrt{2} \cis \left(\dfrac{\pi}{4}\right)$ \\

$z^5 = 4\sqrt{2} \left(\cos \dfrac{\pi}{4} + i \sin \dfrac{\pi}{4}\right) = 4\sqrt{2} \left(\dfrac{\sqrt{2}}{2} + i \dfrac{\sqrt{2}}{2}\right)$ \\

$z^5 = 4 + 4i$ \\

\begin{thebibliography}{999}
%\bibitem{a} Referentie
\end{thebibliography}
\end{document}